\documentclass[titlepage, a4paper, 10pt]{article}

\usepackage{amssymb}
\usepackage{amsmath}
% \usepackage{pgfplots}
\usepackage[margin=1in]{geometry}
\usepackage{fancyhdr}
\usepackage{xcolor}
\usepackage{array}
%\usetikzlibrary{mindmap}

\title{\normalsize \textit{Reduced} PLUS Session Worksheet \\[1em] \Large \textit{MATH 1554:} Linear Algebra}
\date{Summer 2023}
\author{Covering lessons: 2.1 -- 2.5 \& 2.8 -- 2.9}

% ------- Set up header -------
\renewcommand{\sectionmark}[1]{\markright{#1}} % from https://tex.stackexchange.com/a/317393/227665
\pagestyle{fancy}
\fancyhead[L]{Summer 2023}
\fancyhead[C]{\textit{\nouppercase{\rightmark}}}
\fancyhead[R]{\textbf{\thepage}}
% -------------------------------------------------------

% ------- Set up math stuff -------
\renewcommand{\vec}[1]{\mathbf{#1}}
\DeclareMathOperator{\vspan}{Span}
\DeclareMathOperator{\row}{row}
\DeclareMathOperator{\col}{col}
\DeclareMathOperator{\colspace}{Col}
\DeclareMathOperator{\nullspace}{Nul}
\DeclareMathOperator{\rank}{rank}
% -------------------------------------------------------

% ------- Set up Main Idea box environment -------
\newcounter{mainideas}
\setcounter{mainideas}{1}
\newcommand{\mainidea}{\textbf{Main Idea \#\themainideas \stepcounter{mainideas}}}
\setlength\fboxsep{1em}
%\newcommand{\midea}[2]{\begin{center}\fbox{\parbox{#1}{\centering \mainidea #2}} \end{center}} % Could be environment but I couldn't bother

% The way this works is that the savebox collects whatever was between the minipage and the end of it, then in the last line of the definition of the environment
% we regurgigate whatever was collected into an fbox which we center.
\newsavebox{\mybox}
\newenvironment{mideea}[1]
{%
\begin{lrbox}{\mybox}%
	\begin{minipage}{#1\textwidth}
		{\centering \mainidea \par}

}
{\end{minipage}\end{lrbox}\begin{center}\fbox{\usebox{\mybox}}\end{center}}
% -------------------------------------------------------

% ------- Set up activity environment -------
\newcounter{activity}
\setcounter{activity}{1}
\newenvironment{newactivity}
{
	\noindent\textbf{Activity \theactivity}
}
{
\stepcounter{activity}
\par
}
% -------------------------------------------------------

\begin{document}
\maketitle
\part{2.1 -- 2.5}
\section{\centering Matrix Operations (2.1)}
\begin{mideea}{0.9}
	\begin{enumerate}
		\item For \(A \in \mathbb{R}^{m \times n}\) and \(B = \begin{bmatrix}\vec{b}_1 & \vec{b}_2 & \dots & \vec{b}_p\end{bmatrix} \in \mathbb{R}^{n \times p}\), 
	\end{enumerate}
			\footnotesize
                 	\[
                    		   AB = \begin{bmatrix}A\vec{b}_1 & A\vec{b}_2 & \dots & A\vec{b}_p\end{bmatrix}
                         	    	= \begin{bmatrix}\row_1(A) \cdot B \\[0.5em] \row_2(A) \cdot B \\[0.5em] \vdots \\[0.5em] \row_m(A) \cdot B\end{bmatrix}
		         	    	= \begin{bmatrix}
							\row_1(A) \cdot \col_1(B)  & \dots  & \row_1(A) \cdot \col_p(B) \\
							\row_2(A) \cdot \col_1(B)  & \dots  & \row_2(A) \cdot \col_p(B) \\ 
							\vdots                              & \ddots &             \vdots                 \\
							\row_m(A) \cdot \col_1(B) & \dots  & \row_m(A) \cdot \col_p(B)
						\end{bmatrix}
			\]
			\normalsize \centering \textit{Row-column rule for matrix multiplication; equivalent to the other two methods.}
\end{mideea}
	\begin{newactivity}
		Consider the following questions:
		\begin{enumerate}
			\item If \(A \in \mathbb{R}^{3 \times 3}\) and \(B \in \mathbb{R}^{3 \times 3}\), and the third column of \(B\) is the sum of first two columns,
				what can you say about the linear dependency of the columns of \(AB\)? Why is that? (\textit{Hint: \(B = \begin{bmatrix}\vec{b}_1 & \vec{b}_2 & \vec{b}_1 + \vec{b}_2\end{bmatrix}\)}. \(AB = \, ?\))
			\item If the columns of \(A \in \mathbb{R}^{m \times n}\) form a linearly independent set and the columns of \(B \in \mathbb{R}^{n \times p}					\) do not, do the columns of \(AB\) form a linearly dependent set? (\textit{Consider the homogenous equation \(AB\vec{x} = \vec{0}\)}.)
		\end{enumerate}
	\end{newactivity}

\section{\centering Inverses of Matrices (2.2)}
\begin{mideea}{0.9}
	\begin{enumerate}
		\item If \(A \in \mathbb{R}^{n \times n}\) and there exists a matrix \(A^{-1}\) such that \(A A^{-1} = A^{-1}A = I_n\), then \(A\) is invertible.
		\item \(A \ \text{is \emph{non-singular} (invertible)} \iff A\vec{x} = \vec{b} \ \text{is consistent with unique} \ \vec{x} \in \mathbb{R}^n\).
		\item If \(A = \begin{bmatrix}a & b \\ c & d\end{bmatrix}\) and invertible, then 
		\(A^{-1} = \displaystyle \frac{1}{ad -bc} \begin{bmatrix}-a & c \\ b & -d\end{bmatrix} \color{gray} (\text{Note:} \ ab - cd = \det A).\)
	\end{enumerate}
\end{mideea}
\begin{newactivity}
	\begin{enumerate}
		\item What is the inverse of \(A = \begin{bmatrix}1 & 0 & 1 \\ 2 & 1 & 0 \\ 5 & 1 & 1\end{bmatrix}\)? Record the row operations you perform up to 
			the point you reach the REF of \(A\), then complete finding the inverse as usual.
		\item For each recorded row operation in the previous problem, write the corresponding elementary matrix. Multiply these matrices
			together. (This'll help us later on!)
		\item[3.] If \(A \in \mathbb{R}^{n \times n}\) and \(AB = CA = I_n\), is \(B = C\)?
	\end{enumerate}
\end{newactivity}

\newpage
\section{\centering Characteristics of Invertible Matrices (2.3)}
	\begin{mideea}{0.9}
	The Invertible Matrix Theorem (IMT) connects different statements and states that they are equivalent. They concern the characteristics of invertible 
	matrices.
	\end{mideea}
\begin{newactivity}
	Consider the following statements:
	\begin{enumerate}
		\item Is every invertible matrix row equivalent to the identity matrix? Why so?
		\item If a transformation \(T \colon \mathbb{R}^{n} \to \mathbb{R}^n\) defined by \(\vec{x} \mapsto A\vec{x}\) is onto, then \(A\) is 
			invertible.
		\item If a transformation \(T \colon \mathbb{R}^n \to \mathbb{R}^n\) defined by \(\vec{x} \mapsto A\vec{x}\) is one-to-one, then the columns of
			\(A\) necessarily span \(\mathbb{R}^n\).
		\item An invertible matrix can have multiple inverses (e.g. one for each side).
	\end{enumerate}
\end{newactivity}

\section{\centering Partition Matrices (2.4)}
\begin{mideea}{0.9}
	\begin{enumerate}
		\item It is possible to dissect, or partition, a matrix into submatrices. 
		\item All the submatrices in the same column must have the same number of columns. Likewise, all submatrices in the same row must have the 
			same number of rows.
	\end{enumerate}
\end{mideea}

	\begin{newactivity}
		\begin{enumerate}
			\item Given \(B, C, I_n \in \mathbb{R}^{n \times n}\) are all invertible and \(A = \begin{bmatrix}B & I_n \\ 0 & C\end{bmatrix}\), 
				 find \(X, Y, W, Z\) in terms of \(B, C,\) such that \(A^{-1} = \begin{bmatrix}X & Y \\ W & Z\end{bmatrix}\). (\(I_n\) is the identity matrix.)
			\item If \(A = \begin{bmatrix}Q & T \\ H & 0\end{bmatrix},\) where \(Q \in \mathbb{R}^{m \times p}\), \(T \in \mathbb{R}^{q \times l}\), and 
				\(0 \in \mathbb{R}^{q \times l}\), what are the dimensions of \(H\)? Of \(A\)? 
		\end{enumerate}
	\end{newactivity}
\section{\centering \(LU\)-Factorization (2.5)}
\begin{mideea}{0.9}
Many \(m \times n\) matrices could be factored into the product of an \(m \times m\) lower triangular invertible matrix called \(L\) and the matrix's REF, \(U\). To 
find \(L\), keep track of the row-operations that get you to \(U\) and invert and combine them.
\end{mideea}
\begin{newactivity}
	\begin{enumerate}
		% \item If we're working with elementary matrices of dimensions \(n \times n\), what row operation corresponds to the elementary matrices that are their own inverse?
		\item Recall the matrix you found in \textbf{Activity 2} (Question 2). Find the inverse of that matrix, knowing that it's the product of elementary 
			matrices. Then, find the \(LU\)-factorization of \(A = \begin{bmatrix}1 & 0 & 1 \\ 2 & 1 & 0 \\ 5 & 1 & 1\end{bmatrix}\).
		\item Use the \(LU\)-factorization found in the previous question to solve the system \(A\vec{x} = \begin{bmatrix}1 \\ 3 \\ 1\end{bmatrix}\).
	\end{enumerate}
\end{newactivity}
\newpage

\part{2.8 -- 2.9}
\section{\centering Subspaces of \(\mathbb{R}^n\) (2.8)}
\begin{mideea}{0.9}
	\begin{enumerate}
		\item A \textbf{subspace} of \(\mathbb{R}^n\) is a subset of \(\mathbb{R}^n\) that has the following properties:
			\begin{enumerate}
				\item The zero vector is in the subset.
				\item The sum of any two vectors in the subset is in the subset as well.
				\item The multiple of any vector in the subset is in the subset as well.
			\end{enumerate}
		\item For \(A \in \mathbb{R}^{m \times n}\), \(\colspace A\) is the set of all linear combinations of the columns of \(A\). 
		\item For \(A \in \mathbb{R}^{m \times n}\), \(\nullspace A\) is the set of all solutions to the homogenous equation \(A \vec{x} = \vec{0}\).
		\item A \textbf{basis} \(\mathcal{B}\) of a subspace \(H\) is a set of vectors satisfyig two properties:
			\begin{enumerate}
				\item \(\vspan \mathcal{B} = H\).
				\item \(\mathcal{B}\) is a linearly independent set.
			\end{enumerate}
	\end{enumerate}
\end{mideea}

\begin{newactivity}
	\begin{enumerate}
		\item Is every line in \(\mathbb{R}^n\) a subspace?\footnote{When we say ``line", we're talking about a subset of \(\mathbb{R}^n\) where the 
			vectors' ``tip" forms a single point of the line.}
		\item If \(A \vec{x} = A \vec{y} = \vec{b}\), for some \(A \in \mathbb{R}^{m \times n}\) and \(\vec{x}, \vec{y} \in \mathbb{R}^{n}\), is 
			\(2\vec{x} - 2\vec{y} \in \nullspace A\)?
		\item Do the columns of a matrix always form a basis for its column space?
	\end{enumerate}
\end{newactivity}

\section{\centering Dimension \& Rank}
\begin{mideea}{0.9}
	\begin{enumerate}
		\item For a subspace \(H\), \(\dim H = \lvert \mathcal{B}\rvert\) (the numbers of vectors in any of its bases).
		\item \(\dim(\colspace A) = \rank A =\# \ \text{pivot columns in} \ A\).
		\item \(\dim(\nullspace A) = \# \ \text{non-pivot columns in} \ A\).
		\item For any \(A \in \mathbb{R}^{m \times n}\), \(\dim(\colspace A) + \dim(\nullspace A) = \rank A + \dim(\nullspace A) = n.\)
	\end{enumerate}
\end{mideea}
	\begin{newactivity}
		\begin{enumerate}
		\item If \(\mathcal{H} = \{\begin{bmatrix}a + b & a & b\end{bmatrix}^T \in \mathbb{R}^3 \mid a, b \in \mathbb{R}\},\) what is 
			\(\dim \mathcal{H}\)?
		\item If \(A \in \mathbb{R}^{10 \times 16}\), and \(A\) has 10 pivot rows, what is the dimension of the null space of \(A\)?
			\item Find the dimension of the column space and null space of \(A = \begin{bmatrix}1 & 5 & 2 \\ 3 & 15 & 2\end{bmatrix}\).
			\item Find \([\begin{pmatrix}1 & 2 & 1\end{pmatrix}^T]_{\mathcal{B}}\) under the basis \(\mathcal{B}\) that is formed by the
				columns of \(A\) in \textbf{Activity 5} (Question 2).
		\end{enumerate}
	\end{newactivity}

\part{Solutions*}

\section{Part I}
\paragraph{Solution to Activity 1}
\begin{enumerate}
	\item Use the fact that columns of \(AB\) are the linear combinations of the columns of \(A\). If 
		\(B = \begin{bmatrix}\vec{b}_1 & \vec{b}_2 & \vec{b}_1 + \vec{b}_2\end{bmatrix}\). Then, 
		\[AB = \begin{bmatrix}A\vec{b}_1 & A\vec{b}_2 & A(\vec{b}_1 + \vec{b}_2)\end{bmatrix} 
		       = \begin{bmatrix}A\vec{b}_1 & A\vec{b}_2 & A\vec{b}_1 + A\vec{b}_2\end{bmatrix}\]
		Notice that the last column of \(AB\) is the sum of the first two columns, meaning that the columns of \(AB\) are linearly \emph{dependent}.
	\item You could use the above reasoning, but a better way to think about this is going back to the homogenous equation \(AB\vec{y} = \vec{0}\). Note that 
		since the columns of \(B\) are linearly \emph{dependent}, the equation \(B\vec{x} = \vec{0}\) has a nontrivial solution. We could set 
		\(\vec{y} = \vec{x}\) to get \(AB\vec{x} = 0\). But notice that \(AB\vec{x} = A(B\vec{x}) = A\vec{0} = \vec{0}.\) The equation \(AB\vec{x} = \vec{0}\)
		has a nontrivial solution, so its columns are linearly \emph{dependent}. So, the answer is \textbf{yes}.
\end{enumerate}

\paragraph{Solution to Activity 2}
\begin{enumerate}
	\item The first few steps are
		\begin{equation*} 
			\begin{split}
				\small
				&\hphantom{{}\overset{\tiny E_{3}}{\sim}{}}
				\left(\!\!\begin{array}{ccc|ccc}1 & 0 & 1 & 1 & 0 & 0\\ 2 & 1 & 0 & 0 & 1 & 0\\ 5 & 1 & 1 & 0 & 0 & 1\end{array}\!\!\right) 
				\overset{\tiny E_{1}}{\sim}
				\left(\!\!\begin{array}{ccc|ccc}1 & 0 & 1 & 1 & 0 & 0\\ 0 & 1 & -2 & -2 & 1 & 0\\ 5 & 1 & 1 & 0 & 0 & 1\end{array}\!\!\right)
				\overset{\tiny E_{2}}{\sim}
				\left(\!\!\begin{array}{ccc|ccc}1 & 0 & 1 & 1 & 0 & 0\\ 0 & 1 & -2 & -2 & 1 & 0\\ 0 & 1 & -4 & -5 & 0 & 1\end{array}\!\!\right) \\
				&\overset{\tiny E_{3}}{\sim}
				\left(\!\!\begin{array}{ccc|ccc}1 & 0 & 1 & 1 & 0 & 0\\ 0 & 1 & -2 & -2 & 1 & 0\\ 0 & 0 & -2 & -3 & -1 & 1\end{array}\!\!\right) \sim \dots
			\end{split}
		\end{equation*}
		We've reached the REF of the matrix by just applying 3 row operations, which are 
		\[E_1 = \begin{pmatrix}1 & 0 & 0 \\ -2 & 1 & 0 \\ 0 & 0 & 1\end{pmatrix} \quad E_2 = \begin{pmatrix}1 & 0 & 0 \\ 0 & 1 & 0 \\ -5 & 0 & 1\end{pmatrix}
		\quad E_3 = \begin{pmatrix}1 & 0 & 0 \\ 0 & 1 & 0 \\ 0 & -1 & 1 \end{pmatrix}.\]
		Anyways, moving on
		\begin{equation*}
			\begin{split}
			\small
			\dots &\sim
			\left(\!\!\begin{array}{ccc|ccc}1 & 0 & 1 & 1 & 0 & 0\\ 0 & 1 & 0 & 1 & 2 & -1\\ 0 & 0 & -2 & -3 & -1 & 1\end{array}\!\!\right)
			\sim 
			\left(\!\!\begin{array}{ccc|ccc}1 & 0 & 1 & 1 & 0 & 0\\ 0 & 1 & 0 & 1 & 2 & -1\\ 0 & 0 & -2 & -3 & -1 & 1\end{array}\!\!\right) \\
			&\sim 
			\left(\!\!\begin{array}{ccc|ccc}1 & 0 & 0 & -1/2 & -1/2 & 1/2\\ 0 & 1 & 0 & 1 & 2 & -1\\ 0 & 0 & 1 & 3/2 & 1/2 & -1/2\end{array}\!\!\right) = A^{-1}
			\end{split}
		\end{equation*}
	\item We've already written out the elementary matrices in the previous part. Now, to combine them, you \emph{could} multiply them out. A simpler way is 
		to realize that multiplying (lower triangular) elementary matrices essentially ``fills" in the lower part of the matrix. That is,
		\[\begin{pmatrix}1 & 0 & 0 \\ \color{magenta} -2 & 1 & 0 \\ 0 & 0 & 1\end{pmatrix}
		   \begin{pmatrix}1 & 0 & 0 \\ 0 & 1 & 0 \\ \color{teal} -5 & 0 & 1\end{pmatrix}
		   \begin{pmatrix}1 & 0 & 0 \\ 0 & 1 & 0 \\ 0 & \color{red}-1 & 1 \end{pmatrix} 
		= \begin{pmatrix}1 & 0 & 0 \\ \color{magenta}-2 & 1 & 0 \\ \color{teal} -5 & \color{red}-1 & 1 \end{pmatrix}\]

	\item \[{\color{teal} AB} = {\color{orange} CA} = {\color{magenta} I_n} \implies C{\color{teal} AB} = C \cdot {\color{magenta}I_n}\implies ({\color{orange}CA}){\color{teal}B} =  C \implies B = C.\] 
		You could also cite the Invertible Matrix Theorem (IMT), stating that each square matrix has
		only one unique inverse.
\end{enumerate}

\paragraph{Solution to Activity 3}
\begin{enumerate}
	\item \textbf{True} Every invertible matrix \(A \in \mathbb{R}^{m \times n}\) possesses \(n\) pivots, meaning that there is a pivot in every row and column. It follows
		then that the every invertible matrix's RREF is the identity matrix \(I_n\). You could also just cite IMT as the reasoning.
	\item \textbf{True} Again, you can just cite IMT. The main reasoning here is the fact that the mapping \(\vec{x} \mapsto A\vec{x}\) is onto, therefore the standard
		matrix \(A\) must have \(n\) pivots. This means that it's row equivalent to the corresponding idenity matrix \(I_n\), which means it's invertible by IMT.
	\item \textbf{True} Since the transformation is one-to-one, that means that there's \(n\) pivot positions in the matrix. We have \(n\) rows in the matrix (Why?) and
		so the columns of the matrix span \(\mathbb{R}^n\).
	\item \textbf{False} A square invertible matrix cannot have multiple inverses. Look at \textbf{Activity 2} (Question 3).
\end{enumerate}

\paragraph{Solution to Activity 4}
\begin{enumerate}
	\item There are multiple ways to solve this problem. I go with the ``standard" approach. We need a matrix 
		\(A^{-1} = \begin{bmatrix}X & Y \\ W & Z \end{bmatrix}\) such that \(AA^{-1} = A^{-1}A = I_{2n}.\) We need
		\[\begin{bmatrix}B & I_n \\ 0 & C \end{bmatrix}\begin{bmatrix}X & Y \\ W & Z \end{bmatrix} = \begin{bmatrix}I_n & 0 \\ 0 & I_n\end{bmatrix}.\]
		If we multiply the left-side out, we get
		\[\begin{bmatrix}BX + W & BY + Z \\ CW & CZ \end{bmatrix}= \begin{bmatrix}I_n & 0 \\ 0 & I_n\end{bmatrix}.\]
		So, we need
		\begin{align*}
			BX + W &= I \implies X = B^{-1}\\
			BY + Z &= 0 \implies Y = -B^{-1}C^{-1}\\
			CW &= 0 \implies W = 0\\
			CZ &= I \implies Z = C^{-1}
		\end{align*}
	\item The dimensions of \(H\) are \(q \times p\), and the dimensions of \(A\) are \((m + q) \times (p + l)\).

\paragraph{Solution to Activity 5}
\begin{enumerate}
	\item Every invertible matrix could be written as the product of elementary matrices. Since \(B\) is invertible, that means \(B = E_p\dots E_3E_2E_1\). Substituting 
		that in, we get that \(A = (E_p\dots E_3E_2E_1)C\). Since \(C\) is left multiplied by a bunch of elementary matrices, we know that there is a sequence of row
		operations that when applied to \(C\) produce \(A\). Hence, \(A \sim C\).
	\item The matrix we found was \(\begin{pmatrix}1 & 0 & 0 \\ \color{magenta}-2 & 1 & 0 \\ \color{teal} -5 & \color{red}-1 & 1 \end{pmatrix}\). The way 
		inverses work for the product of (lower trangular) elementary matrix is that you negate the entries under the main diagonal. (It might be worthwhile to 
		convince yoursellf why.) Therefore, in this case, 
		\[\begin{pmatrix}1 & 0 & 0 \\ \color{magenta}-2 & 1 & 0 \\ \color{teal} -5 & \color{red}-1 & 1 \end{pmatrix}^{-1} 
		= \begin{pmatrix}1 & 0 & 0 \\ \color{magenta}2 & 1 & 0 \\ \color{teal} 5 & \color{red}1 & 1 \end{pmatrix}.\] The \(LU\)-factorization of the matrix 
		\(A\) is therefore \[\begin{pmatrix}1 & 0 & 1 \\ 2 & 1 & 0 \\ 5 & 1 & 1\end{pmatrix} = \begin{pmatrix}1 & 0 & 0 \\ \color{magenta}2 & 1 & 0 \\ \color{teal} 5 & \color{red}1 & 1 \end{pmatrix}\begin{pmatrix}1 & 0 & 1 \\ 0 & 1 & -2 \\ 0 & 0 & -2 \end{pmatrix}\] 

	\item We need \(\vec{x}\) such that \(A\vec{x} = L\underbrace{(U\vec{x})}_{\vec{y}} = L\vec{y} = \begin{pmatrix}1 \\ 3 \\ 1\end{pmatrix}.\) We can now solve
		for the simpler system \(L\vec{y} = \begin{pmatrix}1 \\ 3 \\ 1\end{pmatrix}\) first then solve for \(U \vec{x} = \vec{y}.\)

		For \(\vec{y}\):
			\[\left(\!\!\begin{array}{ccc|c}1 & 0 & 0 & 1 \\ 2 & 1 & 0 & 3 \\ 5 & 1 & 1 & 1\end{array}\!\!\right) 
			  \sim
			  \left(\!\!\begin{array}{ccc|c}1 & 0 & 0 & 1 \\ 0 & 1 & 0 & 1 \\ 0 & 1 & 1 & -4\end{array}\!\!\right)
			  \sim
			  \left(\!\!\begin{array}{ccc|c}1 & 0 & 0 & 1 \\ 0 & 1 & 0 & 1 \\ 0 & 0 & 1 & -5\end{array}\!\!\right)				
			\]
		So, \(\vec{y} = \begin{pmatrix}1 \\ 1 \\ -5\end{pmatrix}\). Now, to find \(\vec{x}\), solve the system \(U\vec{x} = \vec{y}.\)
			\[\left(\!\!\begin{array}{ccc|c}1 & 0 & 1 & 1 \\ 0 & 1 & -2 & 1 \\ 0 & 0 & -2 & -5 \end{array}\!\!\right) 
			  \sim
			   \left(\!\!\begin{array}{ccc|c}1 & 0 & 1 & 1 \\ 0 & 1 & -2 & 1 \\ 0 & 0 & -2 & -5 \end{array}\!\!\right)
			  \sim \dots \sim
			  \left(\!\!\begin{array}{ccc|c}1 & 0 & 0 & -3/2 \\ 0 & 1 & 0 & 6 \\ 0 & 0 & 1 & 5/2 \end{array}\!\!\right)
			\]
		So, \(\vec{x} = \begin{pmatrix}-3/2 \\ 6 \\ 5/2\end{pmatrix}\).
\end{enumerate}
\end{enumerate}

\section{Part II}

\paragraph{Solution to Activity 6}
\begin{enumerate}
	\item \textbf{No} It's important to realize that by ``line," we mean the subset of vectors of \(\mathbb{R}^n\) such that the ``tip" of the vector touches the line.
				  Anyhow, sometimes we have ``lines" in \(\mathbb{R}^n\) that are not subspaces (such as the solutionset of an inhomogenous system.) The line
				  must pass through the origin in order to be a subspace.
	\item \textbf{Yes} \[A(2\vec{x} - 2\vec{y}) = 2A\vec{x} - 2A\vec{y} = 2\vec{b} - 2\vec{b} = \vec{0}.\]
	\item \textbf{No!}
\end{enumerate}

\paragraph{Solution to Activity 7}
\begin{enumerate}
	\item \textbf{2} Notice that every vector \(\vec{v} \in \mathcal{H}\) could be decomposed as follows:
		\begin{align*}
			\vec{v} = \begin{bmatrix}a + b \\ a \\ b\end{bmatrix} &= \begin{bmatrix}a \\ a \\ 0\end{bmatrix} + \begin{bmatrix}b \\ 0 \\ b\end{bmatrix} \\
											&= a\begin{bmatrix}1 \\ 1 \\ 0\end{bmatrix} + b\begin{bmatrix}1 \\ 0 \\ 1\end{bmatrix}, \quad a, b \in \mathbb{R}
		\end{align*}
		Notice then that \(\mathcal{H}\) is nothing but the span of a linearly independent set of vectors. This set forms a basis for \(\mathcal{H}\). Since there are two
		vectors in the basis, \(\dim \mathcal{H} = 2\).

	\item If there are 10 pivot rows in \(A\), then there are also 10 pivot columns. Hence, \(\dim(\nullspace A) = n - \rank A = 16 - 10 = 6.\)

	\item There are two pivot positions in the matrix. Hence, the dimension of the column space is 2, and the dimension of the null space is 1.
	\item For reference, the matrix is \(\begin{pmatrix}1 & 0 & 1 \\ 2 & 1 & 0 \\ 5 & 1 & 1\end{pmatrix}.\) The basis formed is therefore 
		\[\mathcal{B} = \left\{\!\begin{pmatrix}1 \\ 2 \\ 5\end{pmatrix}, \begin{pmatrix}0 \\ 1 \\ 1\end{pmatrix}, \begin{pmatrix}1 \\ 0 \\ 1\end{pmatrix}\!\right\}.\]
		Since \(\mathcal{B}\) forms a basis, we know that we could write the vector \(\begin{pmatrix}1 \\ 2 \\ 1\end{pmatrix}\) as a linear combination of the basis vectors.
\[\begin{pmatrix}1 \\ 2 \\ 1\end{pmatrix} = c_1\begin{pmatrix}1 \\ 2 \\ 5\end{pmatrix} + c_2\begin{pmatrix}0 \\ 1 \\ 1\end{pmatrix} + c_3\begin{pmatrix}1 \\ 0 \\ 1\end{pmatrix} \implies \begin{pmatrix}1 & 0 & 1 \\ 2 & 1 & 0 \\ 5 & 1 & 1\end{pmatrix}\begin{pmatrix}c_1 \\ c_2 \\ c_3\end{pmatrix} = \begin{pmatrix}1 \\ 2 \\1 \end{pmatrix}\]

		Note \(\left[\!\begin{pmatrix}1 \\ 2 \\1 \end{pmatrix}\!\right]_{\mathcal{B}} = \begin{pmatrix}c_1 \\ c_2 \\ c_3\end{pmatrix}\). Solving the inhomogenous system
		is equivalent to finding the \(\mathcal{B}\)-coordinates of the vector.
		\[\left(\!\!\begin{array}{ccc|c}1 & 0 & 1 & 1\\ 2 & 1 & 0 & 2\\ 5 & 1 & 1 & 1\end{array}\!\!\right) \sim \dots \sim \left(\!\!\begin{array}{ccc|c}1 & 0 & 0 & -1\\ 0 & 1 & 0 & 4\\ 0 & 0 & 1 & 2\end{array}\!\!\right).\]
		So, 
		\(\left[\!\begin{pmatrix}1 \\ 2 \\1 \end{pmatrix}\!\right]_{\mathcal{B}} = \begin{pmatrix}c_1 \\ c_2 \\ c_3\end{pmatrix} = \begin{pmatrix}-1 \\ 4 \\ 2\end{pmatrix}.\)
\end{enumerate}
\end{document}